\documentstyle[12pt,graphicx,amssymb]{article}
%\usepackage{graphicx,amssymb,latexsym}

\title{ An isoperimetric inequality for logarithmic capacity}

\author{ Roger W. Barnard, Kent Pearce and Alexander Yu. Solynin\thanks{This paper 
was finalized during the third author's
visit at the Texas Tech University, Fall 2001. This author thanks
the Department of Mathematics and Statistics of this University for the
wonderful atmosphere and working conditions during his stay in Lubbock.
The research of the third author was supported in part by the
Russian Fund for Fundamental Research, grant no. 00-01-00118a.}}

%\\ \it V.A.Steklov Mathematical Institute,\\ \it Russian Academy of Sciences,\\
%\it Fontanka 27, St.Petersburg,\\ \it 191011, Russia}
%\date{}
%\newcommand{\Bbb}{\bf}

\newcommand{\meas}{{\rm meas}\,}
\newcommand{\area}{{\rm area}\,}
\newcommand{\diam}{{\rm diam}\,}
\newcommand{\CAP}{{\rm cap}\,}

\newcommand{\be}{\begin{equation}}
\newcommand{\ee}{\end{equation}}
\newcommand{\bl}{\begin{lemma}}
\newcommand{\el}{\end{lemma}}
\newtheorem{theorem}{Theorem}
\newtheorem{lemma}{Lemma}
\newtheorem{corollary}{Corollary}
\newtheorem{theo}{Theorem}
\newtheorem{conjecture}{Conjecture}
\newtheorem{lemmaA}{Lemma}
\newcommand{\blemmaA}{\begin{lemmaA}}
\newcommand{\elemmaA}{\end{lemmaA}}
\newtheorem{proposition}{Proposition}
\newcommand{\bprop}{\begin{proposition}}
\newcommand{\eprop}{\end{proposition}}

\newtheorem{fig}{Box for Figure}
\renewcommand{\thetheo}{\Alph{theo}}
\renewcommand{\theequation}{\thesection.\arabic{equation}}
\renewcommand{\thelemmaA}{\Alph{lemmaA}}
%\mag1200
\newcommand{\bt}{\begin{theorem}}
\newcommand{\et}{\end{theorem}}
\newcommand{\bc}{\begin{corollary}}
\newcommand{\ec}{\end{corollary}}
\newcommand{\bcon}{\begin{conjecture}}
\newcommand{\econ}{\end{conjecture}}
\newcommand{\ep}{\varepsilon}
\newcommand{\e}{\varepsilon}

\newcommand{\ph}{\varphi}
\newcommand{\varph}{\varphi}
\newcommand{\C}{{\Bbb {C}}}
\newcommand{\CC}{\overline{\C}}
\newcommand{\U}{{\Bbb{U}}}
\newcommand{\CU}{\overline{\U}}
\newcommand{\th}{\theta}
\newcommand{\var}{\varphi}
\newcommand{\al}{\alpha}
\newcommand{\T}{{\Bbb{T}}}
\newcommand{\N}{{\Bbb{N}}}

\newcommand{\R}{{\Bbb{R}}}
\newcommand{\CR}{\overline{\R}}

\newcommand{\II}{\int\!\!\int}
\begin{document}
\maketitle


\begin{abstract}
We prove a sharp lower bound of the form $\CAP E\ge (1/2)\diam E \cdot \Psi(\area E/((\pi/4) \diam^2 E))$ for the logarithmic capacity 
of a compact connected planar set $E$ in terms of its area and diameter. 
Our lower bound includes as special cases G.~Faber's inequality
$\CAP E\ge \diam E/4$ and G.~P{\' o}lya's inequality $\CAP E\ge (\area E/\pi)^{1/2}$. 
We give explicit formulations, functions of $(1/2)\diam E$, for the extremal domains which we identify.
~\footnote{Key words: logarithmic capacity,  omitted area
problem,  univalent function,  local variation, symmetrization.}
~\footnote{AMS Subject Classification 2000: 30C70, 30C85}

\end{abstract}



\section{Introduction}

The logarithmic capacity, $\CAP E$, of a continuum ( = compact connected set) $E$ in the complex plane $\C$ is defined
by
\be
-\log \CAP E=\lim_{z\to \infty}(g(z)-\log |z|),
\ee
where $g(z)$ denotes Green's function of the unbounded component $\Omega(E)$ 
of $\CC\setminus E$ having singularity at $z=\infty$.

The measure of a set described by the logarithmic capacity is very important in potential theory, analysis,
and PDE's. It combines several characteristics of a compact set, among which are the geometric
concept of transfinite diameter due to M.~Fekete, the concept of Chebyshev's constant from
polynomial approximation, and the concept of the outer radius from conformal mapping, see \cite{D,Du,G,H}.

In general, computation of $\CAP E$ is a difficult problem but there are several estimates of $\CAP E$ in terms
of geometric characteristics of $E$ that are very useful in applications, see \cite{PSz}. For instance,
\be
(1/4)\diam E \le \CAP E\le (1/2)\diam E,
\ee
\be
((1/\pi)\area E)^{1/2}\le \CAP E<\infty.
\ee
The first inequality in (1.2) was found by G.~Faber \cite{F} in a different form. The inequalities in (1.3) and the second inequality in (1.2), 
which is valid for any (not necessary connected) compact set, were proved by G.~P{\' o}lya \cite{Po}. Equality occurs only for rectilinear 
segments in Faber's inequality and for disks in P{\' o}lya's inequalities. The case of equality in the right inequality
in (1.2) was studied by J.~Jenkins \cite{J} and A.~Pfluger \cite{Pf}.

We will employ the following notation throughout this paper:  let $\U =\{z:\,|z|<1\}$ and $\U_r (c) =\{z:\,|z-c|<r\}$, so that $\U = \U_1(0)$.  Finally, let $\U^* = \CC \setminus
\CU$.  

In this paper, we prove the following theorem that contains the left inequalities in (1.2) and (1.3) as special cases.

\bt
Let $E$ be a continuum in $\C$. Then
\be
\CAP E\ge (1/2)\diam E \cdot\Psi(\area E/((\pi/4)\diam^2E)),
\ee
where $1/\Psi(s)$ is a decreasing function from $[0,1]$ onto $[1,2]$  that is the inverse function to $s=p^{-2}[\beta^2(p)-2p(\beta(p)-1)]$, 
with $1\le \beta(p)\le 2$ defined by equation (1.7), with $d$ replaced by $p$.

Equality occurs in (1.4) if and only if $E=a(\C\setminus f_d(\U^*))+b$ for some $a,b\in \C$, $a\not=0$ and $1\le d\le 2$, where the function
$f_d(z)$ is defined for $|z|>1$ by (1.8).
\et


The graph of $\Psi$ is plotted in Figure~1. Figure~2 displays the shape of the extremal continua for $a=1$, $b=0$ and some typical
values of $d$.

Let $s=\area E/((\pi/4)\diam^2E)$. In the case $s=0$, (1.4) gives Faber's inequality and in the case $s=1$, it gives P{\' o}lya's
inequality. Combined with the right-hand side inequality in (1.2) and the classical area-diameter  inequality $0\le \area E/((\pi/4)\diam^2E)\le 1$,
(1.4)  describes the range of one of the quantities $\CAP E$, $\diam E$, or $\area E$ if the other two 
are fixed. Several similar sharp inequalities linking three characteristics of a set are known in geometry. But to
prove such a sharp inequality is a difficult task even for purely geometric quantities, see \cite{JBo}. From this perspective,
(1.4) might be the first known sharp inequality for these three quantities that includes a functional characteristic.



\begin{figure}
$$\includegraphics[scale=0.50]{psi_s_7-n}$$
\caption{Graph of $\Psi(s)$}
\end{figure}


\begin{figure}
$$\includegraphics[scale=.30]{bdy_1p3a-n} \includegraphics[scale=.30]{bdy_1p5a-n} \includegraphics[scale=.30]{bdy_1p9a-n}$$
\caption{Extremal continua}
\end{figure}



Since $\CAP E$, $\diam E$, and $(\area E)^{1/2}$ all change linearly with respect to scaling we can fix one of them,
say $\CAP E$, and then study the region of variability of the other two. In this way, we can reformulate the problem as 
finding the maximal omitted area for the class $\Sigma$ of univalent functions
\be
f(z)=z+a_0(f)+a_1(f) z^{-1}+\ldots
\ee
which are analytic in $\U^*$, except for a simple pole at $z=\infty$. For $ f \in \Sigma$, let $E_f=\C\setminus
f(\U^*)$. 
It is well known, as a consequence of the normalization in (1.5), that for $f \in \Sigma$ that $1 \le (1/2)\diam E_f \le 2$.  
Therefore, for $1\le d\le 2$, we will consider $\Sigma_d=\{f\in \Sigma: \diam E_f=2d\}$. For $f\in \Sigma_d$, define $A_f=\area E_f$ 
and $A(d)=\sup _{f\in \Sigma_d} A_f$. It is well known that 
$$
A_f=\pi(1-\sum_{n=1}^\infty n|a_n(f)|^2)
$$
--- this relation will be often used in what follows.
Theorem~1 is equivalent to 
\bt
Let $f\in \Sigma_d$, $1\le d \le 2$. Then,
\be
A(f)\le \pi[\beta^2-2d(\beta-1)],
\ee
where $1 \le \beta \le 2$  is the unique solution to the equation
\be
d=\beta - (\beta - 1) \log(\beta -1).
\ee



Equality occurs in (1.6) if and only if $f(z)=e^{i\tau}f_d(e^{-i\tau}z)+b$ with $\tau\in \R$, $b\in \C$, and
\be
f_d(z)=d+\int_1^z z^{-1}\varphi(z;d)\,dz,
\ee
where the function $\varphi(z;d)=zf'_d(z)$ is defined by the equation
\be
\varphi(z; d)=A \frac{z^2-1}{z} 
\frac{\sqrt{1+Bz^2+\sqrt{r(z)}}
\sqrt{B+z^2+\sqrt{r(z)}}}
{c(1+z^2)+\sqrt{r(z)}}
\ee
with principal branches of the radicals and
\begin{eqnarray*}
c = \beta - 1\  , \hspace{0.25in} A &=&  \frac{1+c}{2c} \ , \hspace{0.25in} B = 2 c^2 - 1 \  , \\
r(z) &=& 1 + 2Bz^2 + z^4 \ .
\end{eqnarray*}
The function $\varphi$ maps $\U^*$ conformally onto the complement of the ``double anchor"

$$
F(\beta,\psi)=[-i\beta,i\beta] \cup \{\beta e^{it}:\,\frac\pi2-\psi\le t \le \frac\pi2+\psi\}
$$

$$
 \cup \{\beta e^{it}:\,\frac{3\pi}{2}-\psi\le t \le \frac{3\pi}{2}+\psi\},
$$
where $\beta$ is defined by (1.7) and $\psi=(1/2)\cos^{-1}(8\beta^{-1}-8\beta^{-2}-1)$.
\et

\medskip

The graph of the maximal omitted area $A(d)=\pi[\beta^2(d)-2d(\beta(d)-1)]$  is shown in Figure~3. 

\medskip


To prove Theorem~2, we apply techniques developed in \cite{BS}, which were based on symmetrization 
transformations and some elementary  local variations. Section~2 contains preliminary results and necessary definitions. 
In Section~3, we identify the extremal function by solving a specific boundary value problem for analytic
functions. Section~4 completes the proofs of Theorems~1 and 2.

Note that some similar sharp estimates for the area of $f(\U)$ for problems with analytic side conditions
instead of geometric constraints, as imposed in the present paper,  were found in \cite{ASS1,ASS2} using a different method.


\begin{figure}
$$\includegraphics[scale=0.50]{area_d_7-n}$$
\caption{Graph of $A(d)$}
\end{figure}





\section{Preliminaries}
\setcounter{equation}{0}

First we show the existence of an extremal function and describe some simple properties of the maximal omitted area.

\bl
For every $1\le d\le 2$ there is a function $f\in \Sigma_d$ such that $A_f=A(d)$. 

The maximal area $A(d)$ is continuous and strictly decreases from $\pi$ to $0$ as $d$ increases from $1$ to $2$.
\el

{\it Proof}. For a fixed $d$, $\Sigma_d$ is compact in the topology of uniform convergence on compact subsets of $\U^*$. Since
$A_f$ is upper semi-continuous, the existence of an extremal function follows.

Let $1<d_1<d_2\le 2$ and let $f\in \Sigma_d$ be extremal for $A(d_2)$. Note that $f$ has at least one 
non-zero coefficient $a_k(f)$ for some $k\ge 1$. The function
\be
f_t(z)=t^{-1}f(tz)=z+a_0(t)t^{-1}+a_1(t)t^{-2}z^{-1}+\ldots,
\ee
as well as the area $A_{f_t}$ and diameter $d(f_t)=\diam E_{f_t}$ depend continuously on $t$, $1\le t<\infty$. Since $f_t(z)\to z$ as $t\to \infty$,
one can easily show that $d(f_{t})\to 2$ as $t\to \infty$. Hence, there is $t_1>1$ such that $d(f_{t_1})=2d_1$. Therefore
\be
A(d_1)\ge A_{f_{t_1}}=\pi(1-\sum_{n=1}^\infty n t_1^{-2(n+1)}|a_n(f)|^2)>\pi(1-\sum_{n=1}^\infty n|a_n(f)|^2)=A(d_2),
\ee
with strict inequality since $a_k(f)\not=0$. Equation (2.2) proves the strict monotonicity of $A(d)$.

Finally, the compactness of $\Sigma_d$ and continuity of $A_{f_t}$ imply the continuity of $A(d)$.
$\Box$

Since the class $\Sigma_d$ is invariant under the rigid motions of $\C$, i.e. $e^{-i\theta}f(e^{i\theta}z)+b\in \Sigma_d$
if $f\in \Sigma_d$ and $\theta\in \R$, $b\in \C$, we may restrict ourselves to functions $f\in \Sigma_d$ such that the points
$w_1=d$, $w_2=-d$ belong to $E_f$. Thus, the condition $\diam E_f=|w_1-w_2|=2d$ will be assumed if a different condition 
is not imposed explicitly.

To prove symmetry properties of an extremal continuum $E_f$ we shall apply Steiner symmetrization defined as follows:

The Steiner symmetrization of a compact set $E$ w.r.t. the real axis $\R$ is a compact set $E^*$ 
such that for every $u\in \R$, $E^*\cap l(u)=\emptyset$ if $E\cap l(u)=\emptyset$ and $E^*\cap l(u)=\{w=u+it:\,-m\le t\le m\}$ if
$E\cap l(u)\not=\emptyset$. Here $l(u)=\{w=u+it:\,-\infty<t<\infty\}$ and $m=\meas(E\cap l(u))$ denotes the linear Lebesgue
measure. Steiner symmetrization w.r.t. the imaginary or other axis is defined similarly.

It is well known that Steiner symmetrization preserves area and diminishes diameter and logarithmic capacity \cite{H,D}.

\bl
For $1<d<2$, let $f\in \Sigma_d$ be an extremal function normalized as above. Then, $E_f$ possesses Steiner symmetry w.r.t.
the real and imaginary axes.  Moreover,  the boundary of $E_f$, $\partial E_f$, consists of a Jordan rectifiable curve $L_f$ plus, possibly,
some added segments $I_+=[d_0,d]$, $I_-=[-d,-d_0]$, $0<d_0=d_0(d)\le d$, of the real axis.
\el

{\it Proof}. Suppose that $E_f$ does not possess Steiner symmetry w.r.t. $\R$.  Let $E^*$ be the Steiner symmetrization of
$E_f$ w.r.t. $\R$. Note that
\be
2d=\diam E_f=\diam E^* \quad {\mbox{and}}\quad \CAP E_f>\CAP E^*
\ee
since the points $\pm d\in E_f$ and since $E^*$ is not a rigid motion of $E_f$  (see \cite{D}). Let
\be
F(z)=\alpha z+\alpha_0+\alpha_1 z^{-1}+\ldots, \quad \alpha>0,
\ee
map $\U^*$ conformally onto $\Omega(E^*)$. The inequality in (2.3) shows that $\alpha<1$. Let $F_\alpha=\alpha^{-1}F$. Then, 
$F_\alpha\in \Sigma_{d/\alpha}$.   Therefore, we have
$$
A(d/\alpha)\ge A_{F_\alpha}=\pi \alpha^{-2}(1-\sum_{n=1}^\infty n \alpha^{-2}|\alpha_n|^2) \ge \pi (1-\sum_{n=1}^\infty n |a_n(f)|^2)=A(d).
$$
Since $d/\alpha>d$, the latter contradicts the  strict monotonicity of $A(d)$ in Lemma~1.

The same arguments show that $E_f$ possesses Steiner symmetry w.r.t. the imaginary axis.

Let $L^+_f=\{w\in \partial E_f:\,\Im w>0\}$ and $L^-_f=\{w\in \partial E_f:\,\Im w<0\}$. The Steiner symmetries of $E_f$ w.r.t.
the real and imaginary axes can be used to show that $L^+_f$ and $L^-_f$ are rectifiable Jordan arcs; a similar argument was used 
in \cite[Lemma~4]{ASS2}. 
Indeed, let $L^+=\{w\in L^+_f:\,\Re w>0\}$ and let $d_0=\sup \{\Re w:\, w\in L^+ \}$, 
$m_0=\sup \{ \Im w:\,w\in L^+ \}$.  The function
$$
\tau(w)=u+m_0 -v, \quad {\mbox{where}}\quad w=u+iv
$$
is continuous on $ L^+ $ and maps the closure $\bar  L^+ $ one-to-one onto the segment $\{\tau:\,0\le \tau\le d_0+m_0 \}$. Therefore, 
$\bar  L^+ $ is Jordan.  Since $\Re w$ and $\Im w$ are both monotonic on $ L^+ $, it follows that $\bar  L^+ $ is rectifiable. This
implies that $\partial E_f$ consists of a rectifiable Jordan curve $L_f$ plus, possibly, some added  horizontal segments
$[-d,-d_0]$, $[d_0,d]$ and vertical segments $[-im,-im_0]$, $[im_0,im]$ with $0\le m_0\le m<\infty$.

The presence of vertical segments, i.e., the segments  $[-im,-im_0]$, $[im_0,im]$ with $m > m_0$ easily leads to a contradiction: shortening the vertical
slits and expanding the horizontal ones we can find a continuum $\tilde E$ such that 
$\area \tilde E=\area E_f$, $\CAP \tilde E=\CAP E_f=1$, and $\diam \tilde E>\diam E_f$ that contradicts the strict monotonicity
property of $A(d)$ in Lemma~1.
$\Box$

Let $f\in \Sigma _d$ be an extremal function for $A(d)$. By Lemma~2, 
$\partial f(\U^*)=L^+_f\cup L^-_f\cup [-d,-d_0]\cup [d_0,d]$ with $0<d_0\le d$. If $d_0<d$ then 
there is $0<\theta_0<\pi/2$ such that $L^+_f=\{f(z):\, z\in l^+_f\}$, $[d_0,d]=\{f(z):\, z\in l_n^{++}\}$, $[-d,-d_0]=
\{f(z):\,z\in l_n^{+-}\}$, where $l^+_f=\{e^{i\th}:\,\th_0<\th<\pi-\th_0\}$, $l_n^{++}=\{e^{i\th}:\, 0\le \th\le \th_0\}$, 
$l_n^{+-}=\{e^{i\th}:\, \pi-\th_0\le\th\le \pi\}$. The corresponding arcs in the lower half-plane will be denoted by $l_f^-$, $l_n^{-+}$, 
and $l_n^{--}$. The image curves $L_f^+$ and $L_f^-$ are called the {\it free boundary}, the preimages
 $l_f^+$ and $l_f^-$ are called the {\it free arcs}. Respectively, $[d_0,d]$, $[-d,-d_0]$ and $l_n^{++}$, $l_n^{+-}$, $l_n^{-+}$, $l_n^{--}$
are called the {\it non-free boundary} and the {\it non-free arcs}.

To study the behavior of $f'$ on the non-free arcs we shall use two lemmas from \cite{BS}, which are limiting
cases of Theorem~1 in \cite{S2}.

Let $H_\tau^+$ and $H_\tau^-$ be the left and right half-planes w.r.t. the vertical line $l(\tau)=\{w=u+iv:\,u=\tau\}$.
For $D\subset \CC$, let $D_\tau^+=D\cap H_\tau^+$, $D_\tau^-=D\cap H_\tau^-$ and let $D_\tau^*$ denote the set symmetric
to $D$ w.r.t. $l(\tau)$, i.e. $D^*_\tau=\{w=u+iv:\, 2\tau-u+iv\in D\}$.

We say that $D$ possesses the {\it polarization property} in the interval $\tau_1<\tau<\tau_2$ if $(D_\tau^-)_\tau^*\subset D_\tau^+$
for all $\tau_1<\tau<\tau_2$. Extremal configurations shown in Figure~2 give an example of domains possessing the polarization
property in the corresponding intervals $0<\tau<d$.

\bl \cite[Lemma 4]{BS}
Let $f\in \Sigma$, $D=f(\U^*)$, and let $f$ map a boundary arc $\{ e^{i\th}:\, \th_1<\th<\th_2\}$ onto a horizontal interval 
$\{w:\, \Im w=v_0,\tau_1<\Re w<\tau_2\}$. Let $D$ possess the polarization property in $\tau_1<\tau<\tau_2$. Then, 
$|f'(e^{i\th})|$ strictly increases in $\th_1<\th<\th_2$ if $f(e^{i\th_1})=\tau_2+iv_0$ and strictly decreases if $f(e^{i\th_1})=\tau_1+iv_0$.
\el

If $f$ is extremal for $A(d)$, the joint symmetry of $E_f$ assures that the polarization property of $D=f(\U^*)$ holds in
$d_0<\tau<d$. Thus, we obtain from Lemma~3,

\bc
Let $f\in \Sigma_d$ be extremal for $A(d)$, $1<d<2$ and suppose that the non-free arc $l_n^{++}$ is not  degenerate, i.e., $0<\th_0<\pi/2$.
Then, $|f'(e^{i\th})|$ strictly increases in $0<\th<\th_0$ and $\pi<\th <\pi+\th_0$ and strictly decreases in $\pi-\th_0<\th<\pi$ 
and $2\pi - \th_0<\th< 2\pi$.
\ec

We need  a similar result concerning angular polarization. Let $\gamma_\varph=\{w=te^{i\varph}, \, t\ge 0\}$ and let 
$H_\varph^+$ denote the right half-plane w.r.t. the line determined by $\gamma_\varph$ and
$H_\varph^-$ denote the left half-plane w.r.t. the line determined by $\gamma_\varph$.
For $D\subset \CC$, let $D_\varph^+=D\cap H_\varph^+$, $D_\varph^-=D\cap H_\varph^-$ and let $D_\varph^*$ denote the set symmetric
to $D$ w.r.t. the line determined by $\gamma_\varph$ .

We say that a domain $D$ possesses the {\it angular polarization property} in $\varph_1<\varph<\varph_2$ if 
$(D_\varph^-)_\varph^*\subset D_\varph^+$ for all $\var_1<\var<\varph_2$.  For example, domain $G$ depicted in 
Figure~4 possesses the angular polarization property in $0<\varphi <\pi/2$.

\bl \cite[Lemma~5]{BS}
Let $g$ map $\U^*$ conformally onto $D$ and map a boundary arc $\{e^{i\th}:\, \th_1<\th <\th_2\}$ onto a circular arc
$L=\{w=\rho e^{i\varphi}:\, \varphi_1<\varphi<\varphi_2\}$. Let $g(\infty)\in H_{\varphi_1}^+\cap H_{\varphi_2}^+$ and
let $D$ possesses the angular polarization property in $\varphi_1<\varphi<\varphi_2$. Then, $|g'(e^{i\th})|$ strictly
increases in $\th_1<\th<\th_2$ if $g(e^{i\th_1})= \rho e^{i\varphi_1}$ and strictly decreases in $\th_1<\th<\th_2$
if $g(e^{i\th_1})=\rho e^{i\varphi_2}$.
\el


\medskip

{\bf Remark}. The domains $D$ in Lemmas~3 and 4 possess the polarization property for a horizontal interval and 
the angular polarization property w.r.t. 
rays issuing from the origin, resp.  One can easily reformulate these lemmas for arbitrary intervals and for  rays
issuing from  arbitrary centers.  For instance, Lemma~4 in \cite{BS} is formulated for vertical intervals.

The term ``polarization property" comes from the proofs of Theorem~1 in \cite{S2} and Lemmas~4 and 5 in \cite{BS}
that use the polarization transformation.


\medskip

To find a boundary condition for an extremal function $f\in \Sigma_d$ on the free arcs, we apply, in a suitable form, the
local variation used in \cite{BS}. First, we recall the relations linking the logarithmic capacity of a continuum $E$ with
the outer radius and reduced module of $\Omega(E)$. Let
$$
g(w)=w+b_0+b_1 w^{-1}+\ldots
$$
map $\Omega(E)$ conformally onto $\U_R^*=\CC \setminus \CU_R$, where $\U_R=\{\zeta:\,|\zeta|<R\}$. 
The radius $R=R(E)$ of the omitted disk is uniquely determined and is called the outer radius of $\Omega(E)$;
it is well known that $\CAP E=R(E)$, see \cite[\S~10.2]{Du}, \cite[Ch.~7]{G}. The quantity
\be
m(\Omega(E),\infty)=-\frac{1}{2\pi}\log R(E)=-\frac{1}{2\pi}\log \CAP E
\ee
is called the reduced module of $\Omega(E)$ at $w=\infty$. 

The  variation used in \cite{BS} and in the present paper is rather complicated. Let $\Omega \subset \CC$
be a simply connected domain which contains $\infty$ such that its boundary arcs lying in the vicinities of two of its boundary points, $w_1$
and $w_2$, are Jordan and rectifiable. Let $\partial \Omega$ have a tangent $l$ at $w_1$ and let $n_1$ be a unit inward
normal at $w_1$. For $\ep_1>0$ small enough, let $c^0_{\ep_1}$ and $c_{\ep_1}$ be open and closed crosscuts
of $D$ at the boundary point $w_1$, i.e. $c^0_{\ep_1}$ and $c_{\ep_1}$ are respectively the biggest open and closed
arcs of $C_{\ep_1}(w_1)$, where $C_r(w_0)=\{w:\,|w-w_0|=r\}$, such that $w_1+\ep_1n_1\in c^0_{\ep_1}\subset \Omega$
and $w_1+\ep_1n_1\in c_{\ep_1}\subset \bar \Omega$, respectively. Let $\U^+_{\ep_1}(w_1)$ denote the connected component
(half-disk) of $\U_{\ep_1}(w_1)\setminus l$ that contains the point $w_1+\ep_1n_1$ on its boundary. Let $c'_{\ep_1}$ denote the maximal
open circular arc contained in the intersection $c^0_{\ep_1}\cap \partial \U^+_{\ep_1}(w_1)$. Let $\hat \Omega_{\ep_1}$
be a connected component of $\Omega\setminus \overline {\U}_{\ep_1}(w_1)$ containing $\infty$ and let
\be
\Omega_{\ep_1}=\hat \Omega_{\ep_1}\cup \U^+_{\ep_1}(w_1)\cup c'_{\ep_1}.
\ee
Let $I(\ep_1)=\{w=w_1-itn_1:\,-\ep_1<t<\ep_1\}$. For $0<\varphi_1\le 1/2$, let $M(\ep_1,\varphi_1)$ denote the open lune
in $\U_{\ep_1}(w_1)\setminus \U^+_{\ep_1}(w_1)$ bounded by $I(\ep_1)$ and a circular arc $\gamma(\ep_1,\varphi_1)$ that forms angles
of opening $\pi\varphi_1$ with the interval $I(\ep_1)$ at its end points. Let
\be
\Omega(\ep_1,\varphi_1)=\Omega_{\ep_1}\cup M(\ep_1,\varphi_1)\cup I(\ep_1).
\ee

Let $g(w)=g(w;\ep_1,\varphi_1)$ map $\Omega(\ep_1,\varphi_1)$ conformally onto $\U^*$ such that $g(\infty)=\infty$, $g(w_2)=1$. 
Let $0<\varphi_2\le 1/2$ and $\ep_2>0$ be small enough. Let $U^{\ep_2,\varphi_2} \subset \U^*$ be the simply connected domain which 
contains $\infty$ and which is
bounded by the arc $L(\ep_2)=\{e^{i\th}:\, \ep_2 \le |\th| \le \pi \}$ and by the circular arc $L(\ep_2,\varphi_2)$ with ends at the points
$e^{i\ep_2}$ and $e^{-i\ep_2}$ that forms an angle of opening $\pi(1-\varphi_2)$ with the arc $L(\ep_2)$ at the points $e^{i\ep_2}$ 
and $e^{-i\ep_2}$.


The domain
\be
\tilde\Omega=g^{-1}(U^{\ep_2,\varphi_2},\ep_1,\varphi_1)
\ee
will be called  {\it the two point variation of $\Omega$ centered at $w_1$ and $w_2$ with radii $\ep_1$ and $\ep_2$
and inclinations $\varphi_1$ and $\varphi_2$}.

\medskip

The following lemma is a reformulation of Lemma~10 in \cite{BS} for conformal mappings $f$ of $\U^*$ normalized by
condition $f(\infty)=\infty$; in \cite{BS} this result is formulated for conformal mappings $f$ of the unit disk $\U$
with normalization $f(0)=0$.

\bl \cite[Lemma~10]{BS}
Let $w=f(z)$ map $\U^*$ conformally onto $\Omega$ defined above such that $f(\infty)=\infty$, $f(e^{i\th_1})=w_1$, $f(e^{i\th_2})=w_2$
and let there exist the limits
\be
f'(e^{i\th_k})=\lim_{z\to e^{i\th_k},z\in \CU^*} \frac{f(z)-w_k}{z-e^{i\th_k}}\not= 0,\infty \quad {\mbox{for}} \quad k=1,2.
\ee
Let $|f'(e^{i\th_k})|=\al_k$, $k=1,2$. Let $\tilde\Omega(\ep_1,\ep_2,\varphi_1,\varphi_2)$ be the two point variation of $\Omega$ defined
by (2.8) with $\ep_2$ replaced by $\ep_2/\al_2$. Then, for fixed $0 < \varphi_1\le 1/2$ and $0< \varphi_2\le 1/2$,
$$
m(\tilde\Omega(\ep_1,\ep_2,\varphi_1,\varphi_2),\infty)-m(\Omega,\infty)=
$$
\be
\frac{\varphi_1(2+\varphi_1)}{12\pi\al_1^2(1+\varphi_1)^2}\ep_1^2-\frac{\varphi_2(2-\varphi_2)}{12\pi\al_2^2(1-\varphi_2)^2}\ep_2^2+o(\ep_1^2)+o(\ep_2^2)
\ee
and
$$
\area (\C\setminus \tilde\Omega(\ep_1,\ep_2,\varphi_1,\varphi_2))-\area(\C\setminus \Omega)=
$$
\be
-\frac{2\pi\varphi_1-\sin 2\pi\varphi_1}{2\sin^2\pi\varphi_1}\ep_1^2+\frac{2\pi\varphi_2-\sin 2\pi\varphi_2}{2\sin^2\pi\varphi_2}\ep_2^2+o(\ep_1^2)+o(\ep_2^2)
\ee
as $\ep_1\to 0$ and $\ep_2\to 0$.
\el

\medskip

To prove uniqueness of the extremal function in $\Sigma_d$, we need the following modification of Lemma~1 in \cite{S1} where a similar result 
is proved for domains with one axis of symmetry.

\bl (cf. \cite[Lemma~1]{S1}).
For $k=1,2$, let $\Omega_k \subset \CC$ be simply connected domains which contain $\infty$ and which have double symmetry w.r.t. the coordinate axes. 
Let there be a point $\zeta\in \partial \Omega_1$, $\Re \zeta\ge 0$, $\Im \zeta>0$ such that the points $\zeta$, $\bar\zeta$, $-\zeta$, and $-\bar\zeta$
split $\partial \Omega_1$ into four boundary arcs $l_r^+$, $l_r^-$, $l_i^+$, and $l_i^-$, where $l_r^+$ lies in the closed right half-plane
and connects $\zeta$ and $\bar\zeta$, $l_i^+$ lies in the closed upper half-plane and connects $\zeta$ and $-\bar\zeta$. Let
$$
g_k(z)=z+a_1(g_k)z^{-1}+\ldots
$$
map $\U^*$ conformally onto $\Omega_k$.

If $l_r^+\subset \overline{\Omega}_2$, $l_i^+\subset \C\setminus \Omega_2$, then
\be
g_1(r)\le g_2(r), \quad |g_2(ir)|\le |g_1(ir)|
\ee
for all $r>1$. Equality  can occur in (2.12) if and only if $\Omega_1=\Omega_2$.
\el

\section{Boundary value problem for extremal functions}
\setcounter{equation}{0}



In this section $f$ will denote the extremal function in $\Sigma_d$ with $1<d<2$ and $D=f(\U^*)$.

\bl
There is $\beta>0$ such that
\be
|f'(e^{i\th})|=\beta \quad {\mbox{for a.e.}} \quad e^{i\th}\in l_f:=l_f^+\cup l_f^-
\ee
and
\be
|f'(e^{i\th})|<\beta \quad {\mbox{for all}} \quad e^{i\th}\in l_n:=\T\setminus \bar l_f.
\ee
\el

{\it Proof}. Since $\partial D=L_f\cup I_+\cup I_-$ by Lemma~2, where $L_f$ is Jordan and rectifiable, the non-zero finite
limit
\be
f'(\zeta)=\lim_{z\to \zeta, z\in \CU^*}\frac{f(z)-f(\zeta)}{z-\zeta}\not=0,\infty
\ee
exists for almost all $\zeta\in\T$. This easily follows from Theorem~6.8 in \cite{P} applied to the univalent function
$1/f(1/z)$. Assume that
\be
0<\beta_1=|f'(e^{i\th_1})|<|f'(e^{i\th_2})|=\beta_2<\infty
\ee
for some $e^{i\th_1},e^{i\th_2}\in l_f$. Note that (3.3) and (3.4) allow us to apply the two point variation of Lemma~5.

For positive $k_1,k_2$ such that
\be
0<k_1<1<k_2 \quad {\mbox{and }} \quad k_1\beta_1^{-1}>k_2\beta_2^{-1}
\ee
and for fixed $\varphi>0$ small enough consider the two point variation $\tilde D$ of $D$ centered at $w_1=f(e^{i\th_1})$ and
$w_2=f(e^{i\th_2})$ with inclinations $\varphi$ and radii $\ep_1=k_1\ep$, $\ep_2=k_2\ep$, respectively. Computing 
the change in the omitted area by formula (2.11), we find
$$
\area (\C\setminus \tilde D)-\area(\C\setminus D)= \frac{2\pi\varphi-\sin 2\pi\varphi}{2\sin^2\pi\varphi}\ep^2(k_2^2-k_1^2)+o(\ep^2).
$$
Therefore,
\be
\area(\C\setminus\tilde D)>\area(\C\setminus D)
\ee
for all $\ep>0$ small enough. Applying the variation (2.10) of Lemma~6, we get
\begin{eqnarray}
m(\tilde D,\infty)-m(D,\infty)&=&
\frac{1}{12\pi}\left[\frac{\varphi(2+\varphi)}{(1+\varphi)^2}\frac{k_1^2}{\beta_1^2}-\frac{\varphi(2-\varphi)}{(1-\varphi)^2}\frac{k_2^2}{\beta_2^2}\right]\ep^2
+o(\ep^2) \nonumber \\
&=& \left[\frac{\varphi}{6\pi}\left ( \frac{k_1^2}{\beta_1^2}-\frac{k_2^2}{\beta_2^2}\right ) +o(\varphi)\right ]\ep^2+o(\ep^2), 
\end{eqnarray}
which together with (3.5) implies that
\be
m(\tilde D,\infty)>m(D,\infty)
\ee
for all $\ep>0$ small enough if $\varphi$ is chosen such that the expression in the brackets in (3.7) is positive. Let 
$E=\C\setminus D$, $\tilde E=\C\setminus \tilde D$. Equations (3.8) and (2.5) show that
\be
\CAP \tilde E<\CAP E.
\ee
Since $\area \tilde E>\area E$ and $\diam \tilde E\ge \diam E$, (3.9) contradicts the monotonicity property of the function
$A(d)$ in Lemma~1.

\smallskip

Assume that $l_n\not=\emptyset$. Then $f'(1)=f'(-1)=0$. To prove that $|f'(e^{i\th})|<\beta$ for all $e^{i\th}\in l_n\setminus \{\pm 1\}$,
we assume that $\beta=|f'(e^{i\th_1})|<|f'(e^{i\th_2})|=\beta_2$ with $e^{i\th_1}\in l_f$ and some $e^{i\th_2}\in l_n\setminus \{\pm 1\}$.
Then applying the two point variation as above we get (3.6) and (3.9), again contradicting the monotonicity property of $A(d)$.
Hence, $|f'(e^{i\th})|\le \beta$ for all $e^{i\th}\in l_n$, which combined with the strict monotonicity property of Corollary~1
leads to the strict inequality in (3.2).
$\Box$

\bl
If $1<d<2$, then $l_n=\{e^{i\th}:\,-\th_0<\th<\th_0\}\cup \{e^{i\th}:\, \pi-\th_0<\th<\pi+\th_0\}$ with some $0<\th_0=\th_0(d)<\pi/2$; $f'$ is continuous
on $\overline{\U}^*$ and for all $z\in \overline{\U}^*$
\be
|f'(z)|\le |f'(e^{i\th})|=\beta,
\ee
where $e^{i\th}\in \bar l_f$ and $\beta>1.$
\el

{\it Proof}. Consider the function $g(z)=1/f(1/z)$. By Lemma~2, $g$ maps $\U$ onto a Jordan rectifiable domain possibly slit along
two symmetric radial segments lying on the real axis. The double symmetry of $D=f(\U^*)$ implies that $G=g(\U)$ is starlike
w.r.t. $w=0$. Since $G$ is rectifiable and starlike, it follows from classical results of Lavrent'ev, see \cite[p.163]{P}, that $G$ is a Smirnov
domain (non-Jordan in general). This shows that $\log|g'(z)|=\log|f'(1/z)|-2\log|zf(1/z)|$, and therefore $\log|f'(1/z)|$, 
can be represented by the Poisson integral
\be
\log|f'(1/z)|=\frac{1}{2\pi}\int_0^{2\pi}P(r,\th-t)\log|f'(e^{-it})|\,dt
\ee
with boundary values defined a.e. on $\T$, see \cite[p. 155]{P}. Equation (3.11) along with (3.1) and (3.2) shows that $1=|f'(\infty)|\le \beta$ with
 equality only for the function $f(z) \equiv z$.

If $l_n=\emptyset$, then (3.11) and (3.1) show that $|f'|=\beta$ identically on $\U^*$. Therefore, $f(z)\equiv z$ contradicting the 
condition $d=(1/2)\diam \partial f(\U^*)>1$. Hence, $l_n\not=\emptyset$. The latter implies that $f$ is analytic in  vicinities of the points
$z=1$ and $z=-1$ and $f'(z)$ has a simple zero at $z=1$, $z=-1$. Consider the function $h(z)=\log |f'(1/z)/(z^2-1)|$, 
which can be represented by the Poisson integral
\be
h(z)=\frac{1}{2\pi}\int_0^{2\pi}P(r,\th-t)\log|f'(e^{it})/(e^{2it}-1)|\, dt.
\ee
Equation (3.12) and the previous analysis show that $h$ is a bounded harmonic function on $\U$. Let $h_1$ be a bounded harmonic function on $\U$
with boundary values $\log(\beta/|z^2-1|)$ on $l_f$ and $h(z)$ on $l_n$. Then $h_1-h$ has nontangential limit $0$ a.e. on $\T$. Therefore,
$h_1-h \equiv 0$ in $\U$. Hence, $|f'|=\beta$ everywhere on $l_f$.

\smallskip

Since $D$ possesses the double symmetry we need to show only that $f'$ is continuous at $e^{i\th_0}$. By the symmetry principle,
$f$ can be continued analytically through $l_n$ and $f'$ can be continued analytically through $l_f$. This implies that $f$ can 
be considered as a function analytic in a slit disk $\Delta_0=\U_\ep(e^{i\th_0})\setminus [(1-\ep)e^{i\th_0},e^{i\th_0}]$ with
$\ep>0$ small enough.

Using the Julia-Wolff lemma, see \cite[Proposition~4.13]{P}, boundedness of $f'$, and well-known properties of the angular derivatives,
see \cite[Propositions 4.7, 4.9]{P}, one can prove that $f'$ has a finite limit $f'(e^{i\th_0})$, $|f'(e^{i\th_0})|=\beta$, along any
path in $\overline{\U}^*$ ending at $e^{i\th_0}$. The details of this proof are similar to the arguments in Lemma~13 in \cite{BS}.
$\Box$

\bl
Let $f$ be extremal in $\Sigma_d$ for $1<d<2$ and let $\varphi(z)=zf'(z)$. Then $\varphi$ maps $\U^*$ univalently onto
the complement $\Omega(\beta,\psi)=\CC\setminus F(\beta,\psi)$ of the double anchor $F(\beta,\psi)$ defined in Theorem~2 where 
 $\beta$ is defined by $(1.7)$ and $\psi=(1/2)\cos^{-1}(8\beta^{-1}-8\beta^{-2}-1)$.
\el

{\it Proof}. 1) Let $g(z)=f'(\sqrt{z})$. Since $g(\bar z)=\overline{g(z)}$, the symmetry principle implies that $g$ is analytic in $\U^*$.
We will show that $g$ is univalent there.

\begin{figure}
$$\includegraphics[scale=0.40]{g_1p7-n}$$
\caption{Domain $G$ for $\beta=1.7$}
\end{figure}


By Corollary~1, $|g(e^{i\th})|=|f'(e^{i\th/2})|$ strictly increases from $0$ to $\beta$ as $\th$ runs from $0$ to $2\th_0$. Since $\arg
g(e^{i\th})=\arg f'(e^{i\th/2})=(\pi-\th)/2$ strictly decreases from $\pi/2$ to $\varphi_0=\pi/2-\th_0$ as $\th$ runs from $0$ to $2\th_0$,
it follows that $g$ maps the arc $\{e^{i\th}:\, 0\le \th\le 2\th_0\}$ one-to-one onto an analytic Jordan arc $\delta_1$ lying in the domain
$U^+_\beta=\{w\in \U_\beta:\, \Re w>0,\Im w>0\}$ and connecting the points $w=0$ and $w=\beta e^{i\varphi_0}=f'(e^{i\th_0})$.

Since $f(\U^*)$ is starlike w.r.t. $w=0$,
$$
\Re\frac{e^{i\th}f'(e^{i\th})}{f(e^{i\th})}\ge 0
$$
for $0\le \th\le 2\pi$. Since $f(\U^*)$ is symmetric w.r.t. the coordinate axes, the latter inequality shows that
$-\pi \le \arg f'(e^{i\th})\le \pi -\th_0$ for $0\le \th\le \th_0$. This combined with (3.10) implies that $g$ maps the arc
$\{e^{i\th}:\,2\th_0\le \th \le \pi\}$ one-to-one onto the circular arc $\delta_2=\{\beta e^{i\varphi}:\, 0\le \varphi\le \varphi_0\}$
such that $g(e^{2i\th_0})=\beta e^{i\varphi_0}$, $g(-1)=\beta$. By symmetry, $g$ maps the arc $\{e^{i\th}:\,-2\th_0\le
\th \le 0\}$ onto $\bar\delta_1=\{w:\,\bar w\in \delta_1\}$ and the arc $\{e^{i\th}:\,\pi\le \th \le 2\pi-2\th_0\}$ onto $\bar\delta_2=
\{w:\,\bar w\in \delta_2\}$. Thus, $g$ maps the unit circle $\T$ one-to-one onto a closed Jordan arc $\delta$ composed by $\delta_1$,
$\delta_2$, $\bar\delta_2$, and $\bar\delta_1$. Since $g(\infty)=f'(\infty)=1$ the argument principle implies that $g$ maps $\U^*$
conformally and one-to-one onto a simply connected domain $G$ which contains $1$ and which is bounded by $\delta$.  The domain $G$ for $\beta=1.7$ is plotted 
in Figure~4.

The above mentioned properties of $\delta_1$ show that $G$ is circularly symmetric w.r.t. the positive real axis. Therefore by
Lemma~4, $|g'(e^{i\th})|=|f''(e^{i\th/2})| $ strictly decreases in $2\th_0<\th<\pi$.

\medskip

2) Considering boundary values of $\varphi$ we have
$$
\Re \varphi(e^{i\th})=\Re e^{i\th}f'(e^{i\th})=0  \quad {\mbox{for}} \quad 0\le \th \le \th_0
$$
since $\Im f(e^{i\th})=0$ for such $\th$. Since $\Im \varphi(e^{i\th})=|f'(e^{i\th})|$ strictly increases in $0\le \th\le \th_0$, $\varphi$ maps $l_n^{++}$
continuously and one-to-one onto the vertical segment $\{w:\,\Re w=0, 0\le \Im w\le  \beta\}$.

For $\th_0\le \th\le \pi/2$, $|\varphi(e^{i\th})|=\beta$ and 
$$
\frac{\partial}{\partial \th}\arg \varphi(e^{i\th})=
\frac{\partial}{\partial\th}\Im \log (e^{i\th}f'(e^{i\th}))=1+\frac{e^{i\th}f''(e^{i\th})}{f'(e^{i\th})}= 1-\beta^{-1}|f''(e^{i\th})|,
$$
since $e^{i\th}f''(e^{i\th})/f'(e^{i\th})$ is real non-positive for $\th_0\le \th \le \pi/2$.

\medskip

Since $|f''(e^{i\th})|$ strictly decreases in $\th_0\le \th \le \pi/2$ it follows that $\frac{\partial}{\partial \th}\arg \varphi(e^{i\th})$ changes its sign
at most once in the interval $\th_0<\th<\pi/2$.

We have shown in 1) that $\arg f'(e^{i\th})$ decreases from $\pi/2-\th_0$ to $0$  when $\th$ runs from $\th_0$ to $\pi/2$.
Since $\arg \varphi(e^{i\th_0})=\pi/2$, the latter implies that
\be
0<\th_0<\arg \varphi(e^{i\th})=\th+\arg f'(e^{i\th})<\pi-\th_0
\ee
for $\th_0<\th<\pi/2$.

\medskip

We claim that there is $\th_1$, $\th_0<\th_1<\pi/2$ such that 
\be
\begin{array}{ll}
\frac{\partial}{\partial\th}\arg \varphi(e^{i\th})<0 &\quad {\mbox{if}} \quad \th_0<\th<\th_1,\\
\frac{\partial}{\partial\th}\arg \varphi(e^{i\th})>0 &\quad {\mbox{if}} \quad \th_1<\th<\pi/2.
\end{array}
\ee
Suppose to the contrary that  $\frac{\partial}{\partial\th}\arg \varphi(e^{i\th})\le 0$ for all $\th_0<\th<\pi/2$.  Then, we would have that  
$\arg \varphi(e^{i\th})$ monotonically decreases
over $\th_0<\th<\pi/2$. Since $\varphi(e^{i\th_0})=\varphi(i)=i\beta$, we would have
$$
\Delta\,\arg \varphi(e^{i\th})\left|_{\th_0}^{\pi/2}=-2\pi k \quad {\mbox{for some \  }} k\in \N \right.
$$
contradicting (3.13). The assumption $\frac{\partial}{\partial \th} \arg \varphi(e^{i\th})\ge 0$ for all $\th_0<\th<\pi/2$ leads to 
the same contradiction. Since $|f''(e^{i\th})|$ strictly decreases in $\th_0\le \th\le \pi/2$, the claim follows.

\medskip

Let $\psi=\arg \varphi(e^{i\th_1})$. The previous analysis shows that $\th_0<\psi<\pi/2$ and $\varphi$ maps each 
of the arcs $\{e^{i\th}:\,\th_0\le \th\le \th_1\}$ and $\{e^{i\th}:\,\th_1\le \th\le\pi/2\}$ continuously and one-to-one onto
the arc $\{\beta e^{it}:\,\psi\le t\le \pi/2\}$ such that $\varphi(e^{i\th_1})=\beta e^{i\psi}$. The symmetry principle now
yields that $\varphi$ maps the unit circle $\T$ continuously and one-to-one in the sense of boundary correspondence
onto the boundary of the domain $\Omega(\beta,\psi)$. Hence by the argument principle, $\varphi$ maps $\U^*$ conformally
and univalently onto $\Omega(\beta,\psi)$

The normalization $f'(\infty)=1$ leads after some work left to the interested readers to the relation $\psi=(1/2)\cos^{-1}(8\beta^{-1}-8\beta^{-2}-1)$.
$\Box$

\section{Proofs of Theorems 1 and 2}
\setcounter{equation}{0}


To prove uniqueness of the extremal function in $\Sigma_d$, we assume that for some fixed $d$,  $1<d<2$, there are 
distinct extremals $f_1$ and $f_2$. By Lemma~9, $zf'_k(z)=g_{\beta_k}(z)$ for some $1<\beta_1<2$, $1<\beta_2<2$, where $g_{\beta_k}$ maps $\U^*$ conformally onto the domain
$\Omega(\beta_k)=\CC\setminus F(\beta_k,\psi(\beta_k))$.  To be explicit, assume that $\beta_1<\beta_2$. 
The domains $\Omega(\beta_1)$ and $\Omega(\beta_2)$ satisfy the conditions of Lemma~6. Therefore, 
\be
g_{\beta_1}(r)>g_{\beta_2}(r)
\ee
for all $r>1$.

Since $f_k(\U^*)$ is doubly symmetric w.r.t. the real and imaginary axes, it follows that $a_0(f_k) = 0$ for $k=1,2$.  Hence, it follows from the normalization in (1.5) that 
\be
\lim_{R\to \infty}\left ( f_1(R)-f_2(R) \right )= 0.
\ee
On the other hand, since $f_1(1)=f_2(1)=d$, we have
$$
f_1(R)-f_2(R)=\int_1^R t^{-1}(g_{\beta_1}(t)-g_{\beta_1}(t))\,dt
$$
and (4.1) implies that the integrand is positive and hence, that $f_1(R)-f_2(R)$ is a (positive) increasing function of $R$, which 
contradicts  (4.2).

Let $f_d$ denote the  unique extremal function in $\Sigma_d$.
To find $f_d$ explicitly, we represent $\varphi(z;d)=zf'_d(z)$ as
$$
\varphi(z;d)=(F_{d^{-2}}(z^{-2}))^{-1/2},
$$
where $F_p(\zeta) = \zeta+a_2(F_p)\zeta^2+\ldots$, $1/4\le p\le 1$, is the univalent function in the standard class $S$ that maps the unit disk
$\U$ onto the domain $\CC\setminus ((-\infty,-p]\cup \{pe^{i\tau}:\,|\tau-\pi|\le \alpha\})$ with  $\alpha=\cos^{-1}(8\sqrt{p}-8p-1)$. It
is well known that $F_p$ is extremal in a number of problems, for instance in the problem on $\max |a_2(f)|$ 
studied by E.~Netanyahu \cite{N} and T.~Suffridge \cite{S}, 
on the subclass of functions $f\in S$ that cover the disk $\U_p$. Using an explicit expression for $F_p$, see for example,
\cite{S}, we get (1.9) and after an  integration (1.7).

The integral in (1.8) can be evaluated in terms of elementary functions. We leave to the interested readers to check (one can use ``Mathematica" or ``Maple")
that $f_d(1)$ coincides with the right-hand side in (1.7), which in this case is equivalent to the equality $f_d(1)=d$.
Since for each $1\le d\le 2$ the extremal function is unique in $\Sigma_d$, (1.7) has  a unique solution $\beta=\beta(d)$ on the interval $1\le \beta \le 2$;
this also follows easily from the monotonicity of the right-hand side of (1.7).


\medskip


To evaluate the maximal omitted area ${A(d)=\rm{Area}}\,(E_{f_d})$, we apply a standard line integral formula and
the fact that $\Im\, (\bar w\,dw)=0$ on the non-free boundary. We have
$$
A(d)=\frac12 \Im\,\int_{\partial E_{f_d}}\bar w\, dw=\frac12
\Im\,\int_{L_{f_d}}\bar w\,dw=\frac12 \Re\,\int_{l_{f_d}}
\overline{f_d(e^{i\theta})}e^{i\theta}f'_d(e^{i\theta})\,d\theta.
$$
Since $|f'_d|^2=\beta^2$ on $l_{f_d}$, we obtain 
$$
A(d)=\frac{\beta^2}{2}\Re\,\lim_{\e\to +0}\left\{\int_{\e}^{\pi-\e}
\frac{f_d(e^{i\theta})}{e^{i\theta}f'_d(e^{i\theta})}\,d\theta +\int_{\pi +\e}^{2\pi-\e}
\frac{f_d(e^{i\theta})}{e^{i\theta}f'_d(e^{i\theta})}\,d\theta\right\}=
$$
$$
=\frac{ \beta^2}{2} \Im\left\{ \int_{\T} \frac{f_d(z)}{z^2f'_d(z)}\,dz-\pi i {\rm{Res}}\,\left[\frac{f_d(z)}{z^2f'_d(z)},1\right]
-\pi i {\rm{Res}}\,\left[\frac{f_d(z)}{z^2f'_d(z)},-1\right]\right\},
$$
where $\int_{\T}f_d/(z^2f'_d) \,dz$ is understood as the Cauchy principal value. 
The function $f_d/(z^2f'_d)$ has simple poles at $z=1$ and $z=-1$.
Computing the integral and residues, we obtain 
$$
A(d)=\pi[\beta^2-2d(\beta-1)],
$$
which implies (1.6). This finishes the proof of Theorem~2.

\medskip

To deduce (1.4), we write (1.6) in an invariant form:
\be
\frac{\area E}{(\pi/4)\diam^2E}\le p^{-2}[\beta^2(p)-2p(\beta(p)-1)]
\ee
with $p=\diam E/(2\CAP E)$, where $1\le \beta(p)\le 2$ is defined by  (1.7) with $d$ replaced by $p$.
Since the maximal omitted area $A(d)$ strictly decreases, the expression in the brackets in (4.3) decreases and therefore the right-hand side 
of (4.3) itself decreases from $1$ to $0$ when $p$ runs from $1$ to $2$. Therefore there is a function $p=\Psi_1(s)$ inverse to 
$s=p^{-2}[\beta(p)^2-2p(\beta(p)-1)]$. Let $\Psi(s)=1/\Psi_1(s)$.
Since the inverse $\Psi_1$ is decreasing, (4.3) leads to the inequality
$$
p\le \Psi_1(\area E/((\pi/4)\diam^2E),
$$
which is equivalent to (1.4) with equality only for the continua described in Theorem~1.



\begin{thebibliography}{aaaa}


\bibitem[ASS1]{ASS1}
D. Aharonov, H. S. Shapiro and A. Yu. Solynin: A minimal area problem in
conformal mapping, {\it J. Analyse Math}. {\bf 78} (1999), 157--176.

\bibitem[ASS2]{ASS2}
D. Aharonov, H. S. Shapiro and A. Yu. Solynin: A minimal area problem in
conformal mapping II, {\it J. Analyse Math}. {\bf 83} (2001), 259--288.
%{\bf 88} (2001), 157--176.

\bibitem[BS]{BS}
R. W. Barnard and A. Yu. Solynin: Local variations and minimal area problem for Caratheodory
functions, submitted


\bibitem[D]{D}
V. N. Dubinin: Symmetrization in geometric theory of functions of a
complex variable, \emph{Uspehi Mat. Nauk}
 {\bf 49} (1994), 3-76 (in Russian); English translation in \emph{Russian Math. Surveys} {\bf 49}: 1
(1994), 1-79.

\bibitem[Du]{Du}
P.~L.~Duren: Univalent Functions, \emph{Springer-Verlag New York Inc., 1983}.

\bibitem[F]{F}
G. Faber: Neuer Beweis eines Koebe-Bieberbachschen Satzes {\" u}ber konforme Abbildung.  
Sitzgsber.  Math.-phys. Kl. bayer. Akad. Wiss. M{\"u}nchen {\bf 1916}, 39-42.

\bibitem[G]{G}
G.~M.~Goluzin: Geometric theory of functions of a complex variable, \emph{AMS, Providence, R.I. 1969}.

\bibitem[H]{H}
W. K. Hayman: Multivalent Functions, \emph{Cambridge Univ. Press, Cambridge, 1958}.


\bibitem[JBo]{JBo}
I.~M.~Jaglom and V.~G.~Boltjanskii: Convex figures, \emph{Holt, Rinehart and Winston, New York, 1960}.


\bibitem[J]{J}
J. A. Jenkins: A uniqueness result in conformal mapping, \emph{Proc. Amer. Math. Soc.}, {\bf 22} (1969), 324--325.
% Univalent functions and conformal mappings, 2nd ed.
%\emph{Springer, Berlin 1965}.

\bibitem[N]{N}
E.~Netanyahu: On univalent functions in the unit disk whose image contains a given
disk, {\it J. Analyse Math}. {\bf 23} (1970),  305--322.


\bibitem[Pf]{Pf}
A.~Pfluger: On the diameter of planar curves and Fourier coefficients, \emph{Z. Angew. Math. Phys.} {\bf 30} (1979), 305--314.

\bibitem[Po]{Po}
G. P{\' o}lya: Beitrag zur Verallgemeinerung des Verzerrungssatzes auf maehrfach 
zusammenh{\" a}ngende Gebiete, {\it S.-B. Preuss. Akad.} (1928), 280-282.


\bibitem[PSz]{PSz}
G.~P\'{o}lya and G.~Szeg\"{o}: Isoperimetric Inequalities in Mathematical Physics, \emph{Princeton University Press, Princeton, N.J., 1951}.

\bibitem[P]{P}
Ch. Pommerenke: Boundary Behaviour of Conformal Maps,
 \emph{Springer-Verlag, 1992}.

\bibitem[S1]{S1}
A. Yu. Solynin: Some extremal problems in the class $\Sigma(\tau)$, \emph
{Zap. Nauchn. Sem. LOMI} {\bf 196} (1991), 138--153 (in Russian);
English translation in: {\it J. Math. Sci.} {\bf 70} (1994),
no.~6, 2152--2161.


\bibitem[S2]{S2}
A. Yu. Solynin: Functional inequalities via polarization,
{\it Algebra i Analiz} {\bf 8} (1996), 145--185 (in Russian); English
translation in: {\it St. Petersburg Math. J.} {\bf 8} (1997),  1015--1038.


\bibitem[S]{S}
T. Suffridge: A coefficient problem for a class of univalent functions, {\it Mich. Math. J.} {\bf 16} (1969), 33-42.
\end{thebibliography}

\medskip
\noindent
Roger W. Barnard \\
Texas Tech University\\
Lubbock, TX USA\\
email: barnard@math.ttu.edu

\bigskip
\noindent
Kent Pearce \\
Texas Tech University\\
Lubbock, TX USA\\
email: pearce@math.ttu.edu

\bigskip
\noindent
Alexander Yu. Solynin\\
Steklov Institute of Mathematics at St. Petersburg\\
Russian Academy of Sciences\\
Fontanka 27, St.Petersburg\\
191011, Russia\\
and\\
Texas Tech University\\
Lubbock, TX USA\\
email: solynin@pdmi.ras.ru
\end{document}

%\end{document}


Our first variation is elementary.

\bl \cite[Lemma~6]{BS}
For $\ep$ small enough, let $U^\ep=\{z:\,1/\bar z\in \U\setminus \CU_\ep(1-\ep)\}$, where $\U_r(z_0)=
\{z:\,|z-z_0|<r\}$, and let $F^\ep=\partial U^\ep$. Then
\be
\CAP F^\ep=\frac{\pi\ep}{(1-\ep)\sin(\pi\ep/(1-\ep))}=1+\frac{\pi^2}{6}\ep^2+O(\ep^3) \quad {\mbox{as}} \quad \ep\to 0.
\ee
\el


\bibitem[Ba]{Ba}
A. Baernstein $\mathrm{I\!I}$: The size of the set where a univalent function is large, \emph{J. Analyse
Math.} {\bf 70} (1996), 157-173.


\bibitem[B1]{B1}
R. W. Barnard: The omitted area problem for univalent functions. Topics in
complex analysis (Fairfield, Conn., 1983), 53--60, Amer. Math. Soc.,
Providence, R.I., 1985.

\bibitem[B2]{B2}
R. W. Barnard: On open problems and conjectures in complex
analysis. Lecture Notes in Math. {\bf 1435} (1990), 1--26.

\bibitem[BL]{BL}
R. W. Barnard and J. L. Lewis: On the omitted  area problem,
{\it Michigan Math. J.} {\bf 34} (1987), 13--22.

\bibitem[BP]{BP}
R. W. Barnard and K. Pearce: Rounding corners of gearlike domains and
the omitted  area problem,  {\it J. Comput. Appl. Math.} {\bf 14} (1986), 217--226.


\bibitem[Go]{Go}
A. Goodman: Note on regions omitted by univalent functions, \emph{
Bull. Amer. Math. Soc.} {\bf 55} (1949), 363--369.


\bibitem[LSha]{LSha}
M. A. Lavrent'ev and B. V. Shabat:  Methods of the theory of functions of a complex
variable,  4th ed., revised, ``Nauka'', Moscow, 1973; German transl. of 3rd ed.,
VEB Deutscher Verlag Wiss., Berlin, 1967.

\bibitem[L]{L}
J. Lewis: On the minimal area problem, {\it Indiana Univ. Math. J.} {\bf 34}
(1985), 631--661.

\bibitem[M]{M}
M. Marcus: Radial averaging of domains, estimates for Dirichlet integrals and
applications, {\it J. Analyse Math.} {\bf 27} (1974), 47--78.


\bibitem[PrW]{PrW}
M. H. Protter and H. F. Weinberger: Maximum Principles  in
Differential Equations, \emph{Springer-Verlag, New York,  1984}.

\bibitem[S1]{S1}
A. Yu. Solynin: The boundary distortion and extremal problems in
certain classes of univalent functions, \emph
{Zap. Nauchn. Sem. POMI} 204 (1993), 115-142 (in Russian);
English translation in: {\it J. Math. Sci.} {\bf 79} (1996),
no.~5, 1341--1358.


\bibitem[S3]{S3}
A. Yu. Solynin: Modules and extremal metric problems,
{\it Algebra i Analiz} {\bf 11} (1999), no. 1,
3--86 (in Russian); English translation in:
{\it St. Petersburg Math. J.} {\bf 11} (2000), no. 1, 1--65.

\bibitem[S4]{S4}
A. Yu. Solynin: Radial projection and the Poincar\'{e} metric, \emph
{Zap. Nauchn. Sem. POMI} 237 (1997), 148--160;
English translation in: {\it J. Math. Sci.} {\bf 95} (1999),
2267--2275.


\bibitem[We]{We}
A. Weitsman: Symmetrization and the Poincar\'{e} metric,
\emph{Ann. of Math.} {\bf 124}  (1986), 159-169.

\bibitem[W]{W}
V. Wolontis: Properties of conformal invariants, \emph{Amer. J. Math.} 74 (1952), 587-606.